\section{本章纲要：一个帝国如何在不均衡空间中治理}

元的行省、路府、站赤、纸钞和漕运不应被当作统一帝国已经完成的证据。大都、江南、东北、西北、云南与海港具有不同的交通条件、税源、语言和社会组织；国家必须通过官员、书手、军户、商人、寺院和地方中介，才能把命令、粮食、货币与案件接入同一尺度。行省因此是不断调整的治理场域，不是现代省界的预演。

本章将把江南接收、北方供给、海上贸易、宗教机构和元末危机放在同一张关系图中。法典和官文书能够说明制度意图，契约、碑刻、钱币、港口与地方材料则显示局部实践；两者都不能自动替代未留档的人群。读者应持续追问：每项能力由谁承担，其收益流向何处，又在何种区域和情势下失灵。
